\section{The doubled theta product in the original coordinates}

% BEGIN D76K1
\paragraph{D76K1. A rational symplectic change of theta coordinates.}
Use a fixed ordering of $m$ rational positions and an even number
$n=2p$ of rows. Write $T\in\mathfrak H_m$, $Y=\operatorname{Im}T$
and, with elliptic vectors regarded as columns, put
\[
 \mathcal T_f(T,\zeta)=
 \exp\!\left(\frac\pi2\sum_k\zeta_k^{\mathsf t}Y^{-1}\zeta_k\right)
 \sum_x f(x)\exp\!\left(\pi i\sum_k
 (x_k^{\mathsf t}Tx_k+2x_k^{\mathsf t}\zeta_k)\right).
\]
For a rational symplectic map $F$ of period coordinates $(a,b)$
there is a linear isometry $W_F$ on rational Schwartz--Bruhat
tests. Whenever $T'$ is symmetric and a complex matrix $P$ satisfies
$P(a+Tb)=a'+T'b'$ for $(a',b')=F(a,b)$, one has
\[
       \mathcal T_f(T,\zeta)=
       (\det P)^p\mathcal T_{W_Ff}(T',P\zeta).
\]
It also intertwines the Weyl displacement
$U(a,b)f(x)=\exp(\pi i\sum_k(2x_k^{\mathsf t}a_k-a_k^{\mathsf t}b_k))
 f(x-b)$ with $U(F(a,b))$.

Indeed write $F(a,b)=(Aa-Bb,-Ca+Db)$ for the blocks of
$\sigma\in\operatorname{Sp}_{2m}(\mathbb Q)$. Comparing the two
coefficients in the period relation gives
$T'=(AT+B)(CT+D)^{-1}$ and $P=(CT+D)^{-\mathsf t}$ by the
symplectic identities. The denominator is invertible: the graph
of $T$ has positive imaginary pairing and hence intersects the real
vertical coordinate subspace trivially after real symplectic
transport. The same pairing proves $\operatorname{Im}T'>0$.
For $C=0$, take
\[
 (W_F f)(x)=(\det D)^p
 \exp\!\left(-\pi i\sum_k x_k^{\mathsf t}BA^{\mathsf t}x_k\right)
 f(A^{\mathsf t}x).
\]
Here $D=A^{-\mathsf t}$, $BA^{\mathsf t}$ is symmetric, and
the rational change of summation variable cancels this quadratic
with the $B$ term of $T'$. In the Haar norm the finite Jacobian
of $A^{\mathsf t}$ across $n$ rows is $|\det A|^{-n}$ by the
product formula, giving isometry. Expanding the difference of the
quadratic at $x$ and $x-b'$ gives the displacement
$(Aa-Bb,Db)$. For $J=\left(\begin{smallmatrix}0&-I\\ I&0
\end{smallmatrix}\right)$ take $W_F=i^{mp}\widehat f$, with
the positive-sign finite Fourier transform in self-dual measure.
Translation and modulation then interchange by
$(a,b)\mapsto(b,-a)$, including the half-pairing. Poisson summation
on the period cosets of $f$ and the real Gaussian Fourier integral
give the theta factor $(\det T)^p\exp(\pi i\sum_k
\zeta_k^{\mathsf t}T^{-1}\zeta_k)$ for the ordinary theta sum at
$-T^{-1}$. Indeed the Gaussian determinant for $2p$ identical
rows is $(\det(-iT))^{-p}$ and the finite-transform scalar
$i^{mp}$ converts the determinant in the transformed sum to
$(\det T)^p$. The Fourier operator is isometric by finite character
orthogonality. These arguments also preserve finite support modulus
and period. The rational parabolic and $J$ generate: reduce a
Lagrangian against the coordinate Lagrangian by first choosing
a basis of the intersection, then subtract the symmetric graph
on its complementary coordinates by a shear. Exchanging the needed
coordinate pairs restores the coordinate Lagrangian. Each partial
exchange is generated by the two opposite elementary shears
on those coordinates; the lower shear is a conjugate by $J$.
The stabilizer is parabolic. Composition therefore supplies $W_F$
for every $F$. The generator formulas give the ordinary theta
multiplier $(\det(CT+D))^p$ times the exponential with quadratic
$(CT+D)^{-1}C$, equivalently by multiplication of the fractional
transformations. In the Gaussian normalization its quadratic
cancels since
\[
 (P\zeta)^{\mathsf t}(\operatorname{Im}T')^{-1}(P\zeta)
 =\zeta^{\mathsf t}Y^{-1}\zeta
   -2i\zeta^{\mathsf t}(CT+D)^{-1}C\zeta .
\]
This follows by subtracting $T'$ and its conjugate and inverting.
Fourier uniqueness of the ordinary theta coefficients at one fixed
period and all elliptic variables fixes the resulting finite
operator independent of the decomposition.
% END D76K1

% BEGIN D76K2
\paragraph{D76K2. The marked doubled period and the scalar of its product.}
For the native symplectic basis and regular frames $x_1,x_2$ let
$T_i=\tau(x_i)$, $h_i=h(x_i)$ be their Siegel periods and period
normalization matrices. The product of the two framed theta series
at $z,\bar w$, conjugating the second entire value, is
\[
 (\det h)^p\mathcal T_{f_0}(T,h(z,w)),\quad
 T=\operatorname{diag}(T_1,-\bar T_2),\quad
 h=\operatorname{diag}(h_1,\bar h_2),\quad
 f_0(u,v)=f(u)\overline{g(-v)} .
\]
This follows term by term by changing the sign of the second
summation variable. Absolute Gaussian convergence justifies the
product of the two sums.

Define $H=\kappa_0G/2$. For old displacements corresponding to
$y_1,y_2\in E^r$ use the symplectic chart coordinates of $y_1$
in the first block and reverse the sign of the translation
coordinate of $y_2$ in the second block. Their doubled row
coordinates are $d=y_1-y_2$, $b=(y_1+y_2)H/2$. For a rational
basis $e_j$ of $E$ set
\[
 d_a=\sum_j t_{aj}e_j,\qquad
 s_{aj}=\operatorname{Tr}_{E/\mathbb Q}(b_a\bar e_j),\qquad
 F:(\text{old modulation,translation})\longmapsto(s,t).
\]
This map is rational and symplectic. Indeed for any second pair
giving $d',b'$, its alternating pairing is the trace of
$b(d')^*-d(b')^*$. Expanding $d,b$ and using $H^*=-H$ gives
the trace of $y_1H(y'_1)^*-y_2H(y'_2)^*$ (in the trace one
may conjugate the cross terms). In rational chart coordinates
this is the standard symplectic pairing on the first block
minus the one on the second before reversing its translation.
The dimensions agree since $2g=[E:\mathbb Q]r$. Choose one
ordering of the new indices $a,j$ for determinants below.

For precision, write $A_\tau=\mathrm{aLin}_\tau$ and
$B_\tau=\mathrm{aAnti}_\tau$, and denote the two native elliptic
variables by $z_H,z'_H$. First shuffle $(z,w)$ so its first list
takes $z$ in the first $p_0$ positions and $w$ in the others
at each $\tau$; its complementary list reverses the selections.
If these are $u,v$, respectively, the native coordinates are
\[
 z_H=-\tfrac{\tau\kappa_0}{2}G_\tau B_\tau^{-1}u,\qquad
 z'_H=\tfrac{\tau\kappa_0}{2}(A_\tau^{-1}v)^{\mathsf t}G_\tau ,
\]
with $z_H$ a column in this display. The two frame matrices are
invertible: the kernels would identify a vector in the selected
positive space of one frame with a vector in the orthogonal
negative space of the other. Their intersection is zero by strict
definiteness. Equivalently $A_\tau$ takes the positive dual rows
of the second frame and inverse complementary rows of the first;
annihilation forces membership in both incompatible subspaces.
$B_\tau$ is its conjugate with the two frames interchanged.
Thus the two elliptic variables vary independently.

For a Hermitian period $Z$ its rationalization $T'$ is the
symmetric matrix for the quadratic
$2\sum_\tau X_\tau^{\mathsf t}Z_\tau\bar X_\tau$ on rational
coordinates $X_{\tau,a}=\sum_j\tau(e_j)t_{aj}$. Similarly the
rational elliptic vector is
\[
 (L(z,w))_{aj}=\sum_\tau
 \big[\tau(e_j)z_{H,\tau,a}+\overline{\tau(e_j)}z'_{H,\tau,a}\big].
\]
For $Z=\mathrm{dblZ}(x_1,x_2)$ the period relation is
\[
   Lh^{-1}(a+T c)=s+T't,\qquad (s,t)=F(a,c).
\]
Here is a coordinate check. For old periods use respectively
$h_1\operatorname{per}_1(y_1)$ and
$\bar h_2\overline{\operatorname{per}_2(y_2)}$. In place of
$u,v$ the shuffled period entries are
$\tfrac12(B_\tau\overline{y_1+y_2}
       -D_s B_\tau\overline{y_1-y_2})$ and
$\tfrac12(A_\tau(y_1+y_2)+D_s A_\tau(y_1-y_2))$ in embedded
coordinates, where $D_s$ is $-1$ on the first $p_0$ entries and
$+1$ on the rest. Each row selects exactly its defining period.
Insert these expressions in the displayed elliptic formulas. The
primed row is $\tau(b)+\tau(d)Z_\tau$ since
$Z_\tau=(\tau\kappa_0/4)(A_\tau^{-1}D_s A_\tau)^{\mathsf t}G_\tau$.
The other row is $\overline{\tau(b)}+
\overline{\tau(d)} Z_\tau^{\mathsf t}$.
For the latter identity, $B_\tau G_\tau^{-1} A_\tau^{\mathsf t}$
commutes with $D_s$: its two off-diagonal blocks pair a frame
column with a distinct inverse-frame row and vanish. Multiplying
this commutation by the inverses transports $D_s$ between the
two factors. Substitution in $L$ then gives $s$ from the two
$b$ terms by the field trace and $T't$ from the $d$ terms, proving
the period relation.

Write $S$ for the fixed shuffle. The expression for $L$ factors
through $S$, followed by the block inverses of $B_\tau,A_\tau$,
then the two fixed $G_\tau,\tau\kappa_0$ factors and the map of
the embedding basis $e_j$. Thus
\[
 \det L=c\prod_\tau(\det A_\tau\det B_\tau)^{-1}
\]
with a single $c\ne0$ independent of both frames. The embedding
matrix across all embeddings is invertible by separability, and
all other fixed factors are invertible. To compute the Gaussian,
complexify the rational coordinates and change them by this same
embedding matrix to the independent lists $X_\tau,\bar X_\tau$.
The matrix $\operatorname{Im}T'$ in these coordinates has zero
diagonal blocks, and off-diagonal blocks $Y_\tau,Y_\tau^{\mathsf t}$
with $Y_\tau=(Z_\tau-Z_\tau^*)/(2i)$, since its quadratic is
$2\sum X_\tau^{\mathsf t}Y_\tau\bar X_\tau$. Its inverse pairs
the elliptic lists by $2\sum_\tau z'_{H,\tau}Y_\tau^{-1}z_{H,\tau}$.
D76K1 with $P=Lh^{-1}$ cancels the input $\det h$ and contributes
$(\det L)^p$. Consequently for the Schwartz--Bruhat test
$\psi=c^p W_F f_0$ in the rational basis of $E^r$ we obtain
\[
 \Theta_A(f;x_1,z)\overline{\Theta_A(g;x_2,\bar w)}
 = k_0(x_1,x_2;z,w)\,
       \Theta_H(\psi;Z,z_H,z'_H).
\]
Indeed rationalizing the ordinary Hermitian sum gives the rational
Siegel sum at $(T',L(z,w))$ term by term. The residual determinant
product to power $p$ and the exponential of
$\pi\sum_{k,\tau}z'_{H,k}Y_\tau^{-1}z_{H,k}$ are precisely $k_0$.
% END D76K2

% BEGIN D76K3
\paragraph{D76K3. Identification and Heisenberg covariance of the doubled test.}
The $\psi$ of D76K2 is the test chosen by the native definition of
$\mathcal D(f,g)$. In fact $k_0$ is nonzero on frames in the domain.
At fixed frames the two doubled elliptic arguments run through
all values by the invertible maps in D76K2. Vary their real rational
elliptic coordinates over a common period torus for two competing
tests. Fourier integration then recovers each rational coefficient
times its nonzero Gaussian at that period. Thus tests yielding the
product coincide. The constructed expression is linear in $f$ and
conjugate-linear in $g$.

For $y_i$ with chart coordinates $(a_i,c_i)$ the operation on
the second slot is
\[
 \overline{(U(a_2,c_2)g)(-v)}
 =\exp\!\big(\pi i(2v^{\mathsf t}a_2+a_2^{\mathsf t}c_2)\big)
        \overline{g(-v-c_2)}
\]
in one row. Thus the product test is displaced by
$((a_1,a_2),(c_1,-c_2))$. D76K1 carries this to $(s,t)$ of
D76K2. The finite rational pairing of $s$ with the coordinates
of $x$ is $\operatorname{Tr}_{E/\mathbb Q}(bx^*)$. Summing over
rows gives exactly
\[
 \mathcal D(T_{y_1}f,T_{y_2}g)=
    T_H\big(y_1-y_2,(y_1+y_2)\kappa_0G/4\big)\mathcal D(f,g).
\]
Here $T_H(d,b)$ sends $u(x)$ to
$u(x-d)\exp(\pi i\operatorname{Tr}_{E/\mathbb Q}
 \operatorname{tr}(b(2x-d)^*))$. Both the full argument and its
half-translation phase agree with the operator used in L76F4
and O74E2.
% END D76K3
